\section{A geometric evaluation formula that survives \texorpdfstring{$xy=1$}{xy=1}}
\label{cert:sec:spectral}
The preceding sections concern finite identities. We now give an evaluation
method for the full double polylogarithm, including components for which
no finite classical reduction has been established.

\begin{lemma}[One-dimensional integral]
\label{cert:lem:integral}
For $a,b\ge1$, $|x|\le1$, $x\ne1$, and $|y|\le1$,
\begin{equation}
 \Li_{a,b}(x,y)=\frac{x}{(a-1)!}
  \int_0^1\frac{(-\log t)^{a-1}\Li_b(xyt)}{1-xt}\,dt.
\label{cert:eq:integral}
\end{equation}
\end{lemma}
\begin{proof}
For $|x|<1$, the Mellin identity
$n^{-a}=\int_0^1 t^{n-1}(-\log t)^{a-1}\,dt/(a-1)!$
and absolute convergence give
\[
 \sum_{n>m\ge1}(xt)^ny^m/m^b
 =\frac{xt}{1-xt}\Li_b(xyt).
\]
This proves \eqref{cert:eq:integral}. For a nontrivial unit outer argument,
$1-rxt$ stays uniformly away from zero for $r$ sufficiently close to one.
Also $|\Li_b(rxyt)|\le\Li_b(t)$, and
\begin{equation}
 \frac1{(a-1)!}\int_0^1(-\log t)^{a-1}\Li_b(t)\,dt
 =\sum_{m\ge1}\frac1{m^b(m+1)^a}\le\zeta(a+b).
\label{cert:eq:momentbound}
\end{equation}
Dominated convergence and summation by parts in the nested series therefore give the boundary
case. Notice that $xy=1$ causes no problem: the logarithmic endpoint at
$b=1$ is integrable in \eqref{cert:eq:momentbound}.
\end{proof}

Let $T_n$ denote the Chebyshev polynomial defined by
$T_n(\cos u)=\cos(nu)$. Put
\begin{equation}
 d=\sqrt{1-x},\qquad r=\frac{1-d}{1+d},\qquad R=|r|,
\label{cert:eq:dr}
\end{equation}
where the square root has positive real part. For $|x|\le1$, $x\ne1$,
we have $\Real d>0$ and therefore $R<1$.

\begin{lemma}[The kernel expansion]
\label{cert:lem:kernel}
Uniformly for $0\le t\le1$,
\begin{equation}
 \frac1{1-xt}=\frac1d\left(1+2\sum_{n\ge1}r^nT_n(2t-1)\right).
\label{cert:eq:kernel}
\end{equation}
\end{lemma}
\begin{proof}
For $s\in[-1,1]$, summing the two geometric series associated with
$s=\cos u$ gives the classical identity
\[
 1+2\sum_{n\ge1}r^nT_n(s)=\frac{1-r^2}{1-2rs+r^2}.
\]
This is also a direct consequence of the generating function in
\cite{hstruct:DLMF}. Substituting \eqref{cert:eq:dr} shows
\[
 1-xt=d\frac{1-2r(2t-1)+r^2}{1-r^2}.
\]
The series is uniform because $|T_n(s)|\le1$ and $R<1$.
\end{proof}

\begin{definition}
For $v=xy$ and $n\ge0$, define the Chebyshev moment
\begin{equation}
 \mathcal M_n^{a,b}(v)=\frac1{(a-1)!}
 \int_0^1(-\log t)^{a-1}T_n(2t-1)\Li_b(vt)\,dt.
\label{cert:eq:Chebmoment}
\end{equation}
\end{definition}

\begin{theorem}[Certified Chebyshev--moment expansion]
\label{cert:thm:spectral}
Under the hypotheses of Lemma~\ref{cert:lem:integral}, with $x\ne0$,
\begin{equation}
 \Li_{a,b}(x,y)=\frac{x}{d}
 \left(\mathcal M_0^{a,b}(v)+2\sum_{n\ge1}r^n\mathcal M_n^{a,b}(v)\right).
\label{cert:eq:spectral}
\end{equation}
For $K\ge1$, let $E_K$ be the difference between the left side and the
truncation with indices $0\le n<K$. Then
\begin{equation}
 |E_K|\le
 \frac{2|x|\zeta(a+b)}{|d|(1-R)}R^K.
\label{cert:eq:error}
\end{equation}
This estimate is uniform over $|y|\le1$, including $xy=1$.
For nontrivial level-three or level-four outer roots,
\begin{equation}
 |E_K|\le\frac{16}{3\,4^K}.
\label{cert:eq:uniformerror}
\end{equation}
\end{theorem}
\begin{proof}
Insert \eqref{cert:eq:kernel} into \eqref{cert:eq:integral}. The integrable majorant
in \eqref{cert:eq:momentbound} justifies termwise integration and gives
$|\mathcal M_n^{a,b}(v)|\le\zeta(a+b)$ for all $n$.
The tail of the resulting geometric series is \eqref{cert:eq:error}.
For the specified roots, direct substitution in \eqref{cert:eq:dr} gives
$R<1/4$ and $|d|>1$; the code also proves these inequalities by interval
arithmetic before using them. Since $\zeta(a+b)\le\zeta(2)<2$, the
bound becomes $4(1/4)^K/(1-1/4)=16/(3\,4^K)$.
The zero outer argument, omitted only to avoid unnecessary notation,
gives $\Li_{a,b}(0,y)=0$ directly.
\end{proof}

The proof does not appeal to cancellation in a sampled sequence, a fitted
asymptotic expansion, or an empirical convergence diagnostic. In particular,
when the product of two roots is one, the same argument and the same bound
apply. The only new feature is the elementary cancellation in the exact
moment formula, to which we now turn.

\section{Every moment reduces to classical polylogarithms}
\label{cert:sec:moments}
Write the shifted Chebyshev polynomial as
\begin{equation}
 Q_n(t):=T_n(2t-1)=\sum_{j=0}^n c_{n,j}t^j.
\label{cert:eq:Q}
\end{equation}
Its coefficients are integers and can be generated without symbolic
integration:
\begin{equation}
 Q_0=1,\quad Q_1=2t-1,\quad Q_{n+1}=2(2t-1)Q_n-Q_{n-1}.
\label{cert:eq:Qrec}
\end{equation}
Thus
\begin{equation}
 \mathcal M_n^{a,b}(v)=\sum_{j=0}^n c_{n,j}J_{a,b}(j+1;v),
\label{cert:eq:momentpoly}
\end{equation}
where
\begin{equation}
 J_{a,b}(c;v)=\sum_{m\ge1}\frac{v^m}{m^b(m+c)^a},\qquad c\in\Z_{>0}.
\label{cert:eq:J}
\end{equation}
The power series in \eqref{cert:eq:J} is absolutely convergent for $|v|\le1$.

\begin{lemma}[Two-pole partial fractions]
\label{cert:lem:partial}
Put
\begin{align}
 A_s&=(-1)^{b-s}\binom{a+b-s-1}{a-1}c^{-(a+b-s)}
          &&(1\le s\le b),\notag\\
 C_s&=(-1)^b\binom{a+b-s-1}{b-1}c^{-(a+b-s)}
          &&(1\le s\le a).
\label{cert:eq:AC}
\end{align}
Then, for $u\ne0,-c$,
\begin{equation}
 \frac1{u^b(u+c)^a}=\sum_{s=1}^b\frac{A_s}{u^s}
                          +\sum_{s=1}^a\frac{C_s}{(u+c)^s}.
\label{cert:eq:partial}
\end{equation}
In particular, $A_1+C_1=0$.
\end{lemma}
\begin{proof}
The principal part at $u=0$ follows by expanding
$(u+c)^{-a}=c^{-a}(1+u/c)^{-a}$. Its coefficient of $u^{-s}$ after
multiplication by $u^{-b}$ is exactly $A_s$.
At $u=-c$, write $w=u+c$ and expand
$u^{-b}=(-1)^bc^{-b}(1-w/c)^{-b}$; this gives $C_s$.
The difference between the two sides has no finite poles and tends to
zero at infinity, so it is identically zero.
The equality $A_1=-C_1$ follows from
$\binom{a+b-2}{a-1}=\binom{a+b-2}{b-1}$.
\end{proof}

For $s\ge1$, let
$H_c^{(s)}(v)=\sum_{k=1}^c v^k/k^s$ and $H_c^{(s)}=H_c^{(s)}(1)$.
A coefficient outside the ranges in \eqref{cert:eq:AC} is understood to be
zero in the following theorem.

\begin{theorem}[Finite classical moment formula]
\label{cert:thm:moments}
For $0<|v|\le1$, $v\ne1$,
\begin{equation}
 J_{a,b}(c;v)=\sum_{s=1}^b A_s\Li_s(v)
       +v^{-c}\sum_{s=1}^a C_s\bigl(\Li_s(v)-H_c^{(s)}(v)\bigr).
\label{cert:eq:momentfinite}
\end{equation}
At the resonant endpoint $v=1$ the finite expression is instead
\begin{equation}
 J_{a,b}(c;1)
 =-C_1H_c^{(1)}+
   \sum_{s=2}^{\max(a,b)}(A_s+C_s)\zeta(s)
   -\sum_{s=2}^a C_sH_c^{(s)}.
\label{cert:eq:momentone}
\end{equation}
No regularized zeta value is needed. At $v=0$, $J_{a,b}(c;0)=0$.
\end{theorem}
\begin{proof}
Multiply \eqref{cert:eq:partial} by $v^m$ and sum over $m\ge1$.
For $|v|<1$, all rearrangements are absolute, and
\[
 \sum_{m\ge1}\frac{v^m}{(m+c)^s}
 =v^{-c}\bigl(\Li_s(v)-H_c^{(s)}(v)\bigr).
\]
Radial limits give \eqref{cert:eq:momentfinite} at $|v|=1$, $v\ne1$.
At $v=1$, combine the two $s=1$ pieces at the level of partial sums:
\[
 A_1\sum_{m=1}^N\frac1m+C_1\sum_{m=1}^N\frac1{m+c}
 =-C_1H_c^{(1)}+C_1\sum_{m=N+1}^{N+c}\frac1m.
\]
The last term tends to zero. Every $s\ge2$ term converges separately,
which gives \eqref{cert:eq:momentone}.
\end{proof}

For example, when $a=b=1$, \eqref{cert:eq:momentone} reduces to
\begin{equation}
 J_{1,1}(c;1)=\frac{H_c}{c}.
\label{cert:eq:J11}
\end{equation}
Consequently the resonant depth-two function has the explicit expansion
\begin{equation}
 \Li_{1,1}(x,x^{-1})=
 \frac{x}{\sqrt{1-x}}\left[
 1+2\sum_{n\ge1}r^n\sum_{j=0}^n c_{n,j}\frac{H_{j+1}}{j+1}
 \right],
\label{cert:eq:resonant11}
\end{equation}
valid for a nontrivial unit outer argument. Its coefficients are rational.
The error after $K$ terms is bounded by \eqref{cert:eq:error}, and the uniform
level-three/four bound \eqref{cert:eq:uniformerror} applies.
This is an exact infinite identity for a conditionally
convergent depth-two case; it is not a finite rational expression in
classical constants.

\subsection{A resonant example beyond the first weight}
For $(a,b)=(2,1)$, the coefficients are
$A_1=c^{-2}$, $C_1=-c^{-2}$, and $C_2=-c^{-1}$, giving
\begin{equation}
 J_{2,1}(c;1)=\frac{H_c}{c^2}+\frac{H_c^{(2)}-\zeta(2)}{c}.
\label{cert:eq:J21}
\end{equation}
Substitution in \eqref{cert:eq:momentpoly} and \eqref{cert:eq:spectral} yields a
geometrically convergent formula for
$\Li_{2,1}(\rho,\rho^{-1})$, $\rho=e^{2\pi\iu/3}$,
whose coefficient data involve only rational numbers and $\zeta(2)$.
The location of the singular product $v=1$ has become an explicit,
finite harmonic-sum correction rather than a numerical failure mode.

\section{Optimal geometric rate and the conductor dependence}
\label{cert:sec:optimal}
The rate in Theorem~\ref{cert:thm:spectral} is not merely an artifact of an
arbitrary expansion point. It is the best exponential rate obtainable
by uniformly approximating its outer rational kernel with polynomials.
This statement concerns the kernel, not the minimal complexity of all
possible algorithms for evaluating a special value.

\begin{theorem}[Sharp exponential rate for polynomial kernels]
\label{cert:thm:optimal}
Fix $0<|x|\le1$, $x\ne1$, and let $d,r,R$ be as in \eqref{cert:eq:dr}.
For $K\ge1$, define
\[
 e_{K-1}(x)=\inf_{\deg p<K}\max_{0\le t\le1}
           \left|\frac1{1-xt}-p(t)\right|,
\]
where the infimum allows complex polynomial coefficients. Then
\begin{equation}
 \frac{R^K}{|d|}\le e_{K-1}(x)
       \le\frac{2R^K}{|d|(1-R)}.
\label{cert:eq:bestpoly}
\end{equation}
In particular,
$\lim_{K\to\infty}e_{K-1}(x)^{1/K}=R$.
\end{theorem}
\begin{proof}
The upper bound is the uniform tail bound in \eqref{cert:eq:kernel}.
For the lower bound, use the $K$th Chebyshev coefficient of
$f(s)=(1-x(s+1)/2)^{-1}$. By \eqref{cert:eq:kernel}, it is $2r^K/d$.
The corresponding coefficient of a polynomial of degree less than $K$
vanishes. For any continuous complex function $g$ on $[-1,1]$,
\[
 \left|\frac2\pi\int_0^\pi g(\cos u)\cos(Ku)\,du\right|
 \le2\|g\|_\infty.
\]
Apply this with $g=f-p$ to obtain
$2R^K/|d|\le2\|f-p\|_\infty$ and take the infimum.
Taking $K$th roots proves the final statement.
\end{proof}

\begin{corollary}[Adjacent roots]
\label{cert:cor:conductor}
For $x_q=e^{2\pi\iu/q}$, as $q\to\infty$,
\begin{equation}
 -\log R_q=2\sqrt{\frac\pi q}+O(q^{-3/2}).
\label{cert:eq:rateq}
\end{equation}
For fixed $a,b$, the bound \eqref{cert:eq:error} guarantees an absolute error
at most $10^{-P}$ using
\begin{equation}
 K=O\bigl(\sqrt q(P+\log q)\bigr)
\label{cert:eq:Kq}
\end{equation}
outer-kernel terms, uniformly in $|y|\le1$.
\end{corollary}
\begin{proof}
With $\theta=2\pi/q$,
$d=\sqrt{1-e^{\iu\theta}}
=e^{-\iu\pi/4}\sqrt\theta\,(1+O(\theta))$.
Hence
\[
 -\log|r|=\Real\left(\log(1+d)-\log(1-d)\right)
 =2\Real(d+d^3/3+\cdots)
 =\sqrt{2\theta}+O(\theta^{3/2}),
\]
which is \eqref{cert:eq:rateq}. Also $|d|\asymp q^{-1/2}$ and
$1-R\asymp q^{-1/2}$, so the prefactor in \eqref{cert:eq:error} is $O(q)$.
Solving the resulting geometric inequality gives \eqref{cert:eq:Kq}.
\end{proof}

For comparison, expansion around $t=1/2$ gives
\[
 \frac1{1-xt}=\frac1{1-x/2}
 \sum_{n\ge0}\left(\frac{x(t-1/2)}{1-x/2}\right)^n.
\]
At a unit outer argument its uniform ratio is
$R_{\mathrm{mid}}=1/|2-x|$. For $x=x_q$,
\begin{equation}
 -\log R_{\mathrm{mid}}
 =\tfrac12\log(5-4\cos(2\pi/q))
 =\frac{4\pi^2}{q^2}+O(q^{-4}).
\label{cert:eq:midrate}
\end{equation}
Thus its corresponding sufficient degree is
$O(q^2(P+\log q))$. This is a comparison between two proved expansions,
not a claim to improve every previously published polylogarithm algorithm.

\paragraph{Arithmetic cost versus approximation degree.}
With weights fixed, the monomial moments for $1\le c\le K$ can be
assembled from finitely many single values and incrementally updated
$H_c^{(s)}(v)$ in $O((a+b)K)$ arithmetic operations.
The triangular transform \eqref{cert:eq:momentpoly} costs $O(K^2)$ operations
in the direct implementation. Thus, after the necessary classical values
are available, $O(K^2+(a+b)K)$ arithmetic operations suffice for the
spectral truncation. These are unit-cost arithmetic counts, not bit
complexity bounds. At general large conductor, computing the inner single
values and representing the algebraic arguments have their own costs.
The implementation supplied here certifies levels three and four; the
analytic theorem has the wider stated domain.

\section{Outward-rounded evaluation and the treatment of cancellation}
\label{cert:sec:cert}
A theorem about a truncation error is not by itself a certificate for a
floating-point implementation. This section specifies how finite
arithmetic errors are enclosed in the accompanying executable code.

\subsection{An independent geometric expansion for the single values}
To avoid using the desired parity identities inside their own numerical
checks, the certified evaluator computes its classical constants from a
separate Euler expansion.

\begin{lemma}[Euler expansion with a simple rigorous tail]
\label{cert:lem:euler}
Let $s\ge1$ be an integer and let $z\ne1$, with
$\tau=z/(z-1)$ satisfying $|\tau|<1$. For the arguments used below,
\begin{equation}
 \Li_s(z)=-\sum_{N\ge1}\tau^N\frac{h_{s-1}^{\star}(N)}{N},
\label{cert:eq:euler}
\end{equation}
where $h_d^{\star}(N)$ is the complete homogeneous symmetric polynomial
of degree $d$ in $1,1/2,\ldots,1/N$, and $h_0^{\star}(N)=1$.
If $|\tau|\le r_0<1$, then the error after $N\le M$ is at most
\begin{equation}
 \frac{r_0^{M+1}}{1-r_0}.
\label{cert:eq:eulertail}
\end{equation}
\end{lemma}
\begin{proof}
In the standard Mellin integral for $\Li_s(z)$, expand
\[
 \frac{z}{1-zt}=-\sum_{N\ge1}\tau^N(1-t)^{N-1}.
\]
The coefficient integral is
\[
 \frac1{(s-1)!}\int_0^1(-\log t)^{s-1}(1-t)^{N-1}\,dt
 =\frac{h_{s-1}^{\star}(N)}{N}.
\]
Indeed, summing its generating series in $u^{s-1}$ gives
\[
 \int_0^1 t^{-u}(1-t)^{N-1}\,dt
 =\frac1N\prod_{j=1}^N(1-u/j)^{-1},
\]
first for $u$ close to zero. This can be proved by the beta integral or
by $N-1$ integrations by parts. Each coefficient integral is between
zero and one, since $(1-t)^{N-1}\le1$. The geometric tail then proves
\eqref{cert:eq:eulertail}.
\end{proof}

For the nontrivial roots of levels three and four, and for $z=-1$,
$|\tau|\le1/\sqrt2<3/4$.
The implementation uses the rational choice $r_0=3/4$ and computes the
coefficients incrementally from
\begin{equation}
 h_d^{\star}(N)=h_d^{\star}(N-1)+\frac{h_{d-1}^{\star}(N)}{N}.
\label{cert:eq:hrec}
\end{equation}
It obtains $\zeta(s)$ for $s\ge2$ from
$-\Li_s(-1)/(1-2^{1-s})$, obtains $\pi$ from
$4\Imag\Li_1(\iu)$, and obtains $L$ from $-\Li_1(-1)$.
No conjectural reduction enters these evaluations.

\subsection{The interval arithmetic contract}
Fix a decimal scale $Q=10^D$. A real interval is stored as integers
$[l,u]/Q$. Addition is exact at this scale. For multiplication, the four
integer endpoint products are enclosed by a floor of their minimum
and a ceiling of their maximum after division by $Q$.
A reciprocal is permitted only when zero is excluded, and is again
rounded outward. A nonnegative square root is enclosed with integer
square roots of $lQ$ and $uQ$. Complex values are rectangular pairs of
such real intervals.

All these rules are elementary rational inequalities. By induction over
the expression tree, the computed rectangle contains the exact finite
expression. The analytic disk radius from \eqref{cert:eq:error} or
\eqref{cert:eq:eulertail} is then added to both real and imaginary coordinate
intervals. Enclosing a disk by a square is conservative but valid.
The square-root branch in \eqref{cert:eq:dr} is fixed by its positive real
part and the sign of the imaginary part of $1-x$.

The executable checks $R<1/4$ and $|d|\ge1$ by interval inequalities
before using \eqref{cert:eq:uniformerror}. It rejects unsupported outer
arguments instead of silently using an unproved bound. In particular,
the supplied certified implementation does not accept outer argument one.

\subsection{Why extra working precision is necessary}
The coefficients in \eqref{cert:eq:Q} alternate in sign and satisfy
\begin{equation}
 \sum_{j=0}^n|c_{n,j}|=T_n(3)
 =\frac{(3+2\sqrt2)^n+(3-2\sqrt2)^n}{2}.
\label{cert:eq:cnnorm}
\end{equation}
The alternating signs follow, for example, from the explicit formula
\[
 c_{n,j}=(-1)^{n-j}\frac{n}{n+j}\binom{n+j}{2j}4^j
 \quad(n\ge1),
\]
and evaluating $Q_n$ at $t=-1$ gives \eqref{cert:eq:cnnorm}.
Thus computing the moments by monomial conversion can lose a number of
digits proportional to $K$, even though the moments themselves are bounded.

The interval code does not assume that this cancellation is harmless.
For target $P=35$ it uses $K=73$ spectral terms and $D=246$ working decimal
places. Each single value uses $M=1997$ Euler terms and its own rigorous
tail. These are conservative choices rather than tuned speed records.
The actual final widths, not a heuristic guard-digit count, determine
whether a certificate is accepted. A future stable recurrence or validated
Clenshaw-style evaluation could reduce the guard overhead; that is not
needed for the present certificates.
